\chapter{Introduction}\label{ch:intro}

\section{Background}
Modern hospitality e-commerce and online travel platforms rely heavily on recommendation engines to assist users in navigating extensive options. In destinations such as Delhi NCR, where travelers evaluate hundreds of accommodation properties across diverse budgets, locations, and amenities, effective personalization is critical. Traditional recommendation platforms typically compute user-item affinities using collaborative filtering (CF) algorithms, matrix factorization (such as Singular Value Decomposition), or content-based similarity models.

However, conventional recommendation systems suffer from significant usability and trust limitations. Most algorithms operate as black boxes, presenting users with ranked lists of items without transparent justifications for *why* a specific hotel was selected for a specific traveler profile. Furthermore, while recent advances in conversational Artificial Intelligence and Large Language Models (LLMs) enable intuitive natural language query interfaces, deploying naive Retrieval-Augmented Generation (RAG) architectures introduces severe reliability risks. Large Language Models frequently generate plausible yet unsupported claims—commonly termed hallucinations—such as fabricating non-existent hotel amenities or misrepresenting guest review sentiment.

To address these challenges, modern recommendation platforms require a unified architecture that integrates preference-aware ranking, multi-dimensional aspect explainability, review-grounded text retrieval, automated hallucination protection, and an enterprise database infrastructure capable of enforcing auditable data lineage.

\section{Problem Statement}
The core problem addressed by this project is the lack of explainability, grounding, and infrastructure auditability in automated travel recommendation systems. Specifically, existing platforms exhibit four key deficiencies:
\begin{enumerate}
    \item \textbf{Opaque Ranking Logic}: Standard recommendation models generate scalar prediction scores without exposing how individual property attributes (such as cleanliness, service, or value for money) align with user preferences.
    \item \textbf{LLM Hallucination and Uncited Responses}: Conversational travel search engines frequently output ungrounded assertions regarding hotel features due to a lack of strict context compression, inline citation enforcement, and post-generation validation.
    \item \textbf{Data Drift and Fragile Vector Management}: File-based vector databases and raw dataset files operating independently of relational metadata suffer from data drift, non-deterministic similarity searches, and a lack of ACID transaction guarantees.
    \item \textbf{Uncoordinated Data Ingestion}: Upstream data collection, text cleaning, sentiment extraction, and database loading scripts executed in isolation create operational fragmentation, increasing the risk of unvalidated database overwrites and redundant embedding recalculations.
\end{enumerate}

Operationally, \textit{TrustLayer-AI} defines \textbf{trust} as the verifiable alignment between explicit user preference constraints, quantitative aspect-level sentiment scores, retrieved review text evidence, and deterministic database state.

\section{Motivation}
Developing a robust recommendation platform requires extending beyond basic model fitting to address real-world system evolution challenges. During the iterative development of TrustLayer-AI, initial collaborative filtering experiments revealed severe performance degradation ($\text{NDCG}@10 = 0.006$) caused by a 99.27\% user-item matrix sparsity. Attempting to combine collaborative filtering with content-based models using linear score blending failed because cosine similarity scores (concentrated between 0.8 and 0.9) and predicted rating scores (ranging from 1.0 to 5.0) exhibited severe scale mismatches, causing optimization algorithms to completely disable the content-based component. Resolving this required implementing Reciprocal Rank Fusion (RRF), which combines ordinal ranks rather than uncalibrated raw scores.

Similarly, providing user-facing explanations using SHAP (SHapley Additive exPlanations) approximations proved computationally unfeasible due to high inference latency. This motivated the design of an analytical feature-matching explainer engine that computes direct aspect alignment scores in real time.

Finally, transitioning from legacy file-based stores (CSV files and ChromaDB) to an enterprise PostgreSQL 17.6 database engine equipped with the \texttt{pgvector} extension was motivated by the need for unified relational and vector ACID transactions. Managing this transition cleanly required building a 9-stage repeatable data ingestion engine and a single-command master CLI orchestrator equipped with live terminal progress monitoring and signal safety.

\section{Research and Engineering Questions}
This Independent Project is guided by five primary research and engineering questions:
\begin{itemize}
    \item \textbf{EQ1 (Hybrid Ranking Fusion)}: How can collaborative filtering and content-based feature similarities be effectively combined using Reciprocal Rank Fusion (RRF) to resolve score calibration mismatches and mitigate extreme matrix sparsity?
    \item \textbf{EQ2 (Real-Time Explainability)}: How can multi-dimensional Aspect-Based Sentiment Analysis (ABSA) scores be structured to generate deterministic, real-time feature-alignment explanations without incurring the latency overhead of SHAP approximations?
    \item \textbf{EQ3 (Grounded Retrieval and Hallucination Protection)}: How can hybrid retrieval (combining vector similarity, metadata filtering, and recommender reranking) be coupled with token-budget context compression and post-generation grounding validation to achieve $>95\%$ grounded response rates and $<2\%$ hallucination rates?
    \item \textbf{EQ4 (Unified Storage Parity)}: How can legacy CSV files and vector databases be migrated to a unified PostgreSQL 17.6 + \texttt{pgvector} storage engine while maintaining 100\% embedding cosine similarity and recommendation parity?
    \item \textbf{EQ5 (Auditable Pipeline Orchestration)}: How can a multi-stage data engineering pipeline enforce SHA-256 field-level diffing, dry-run safety, human approval gates, selective vector synchronization, and live terminal progress tracking with signal protection?
\end{itemize}

\section{Objectives}
To answer these engineering questions, the project fulfilled the following objectives:
\begin{enumerate}
    \item \textbf{Data Engineering and NLP Sentiment Extraction}: Acquire and clean hotel metadata and review text for 1,661 properties across Delhi NCR; deploy DistilBERT models to extract sentence sentiment probabilities and compute 5 aspect sentiment scores (Cleanliness, Service, Location, Value, Staff Behavior).
    \item \textbf{Hybrid Recommendation Modeling}: Build Content-Based filtering, Collaborative Filtering (SVD), and Reciprocal Rank Fusion (RRF, $k=60$) models to generate personalized hotel rankings.
    \item \textbf{Analytical Explainability Layer}: Develop an analytical feature-matching explainer engine (`\texttt{explainer.py}`) to output transparent aspect alignment scores and qualitative explanation badges.
    \item \textbf{Grounded RAG and Hallucination Protection}: Segment review text into 7,910 chunks, construct a Hybrid Retrieval engine, and build a grounding validation pipeline (`\texttt{grounding\_validator.py}`) that intercepts and strips unverified amenity claims from LLM responses.
    \item \textbf{PostgreSQL 17.6 and pgvector Migration}: Unified relational metadata and 384-dimensional vector embeddings (`\texttt{all-MiniLM-L6-v2}`) inside PostgreSQL 17.6, establishing zero orphan records across an 18-test master provenance suite.
    \item \textbf{Repeatable Ingestion and CLI Orchestration}: Build a 9-stage repeatable data ingestion engine featuring SHA-256 field-level diffing, dry-run safety (`\texttt{dry\_run.json}`), human approval gates, selective vector synchronization, and a single-command master CLI orchestrator (`\texttt{scripts/orchestrator.py}`) with real-time ASCII progress tracking and Ctrl+C interrupt safety.
\end{enumerate}

\section{Scope of the Project}
The scope of this project encompasses:
\begin{itemize}
    \item \textbf{Geographic Domain}: Accommodation properties located within the Delhi NCR region (1,661 canonical hotels).
    \item \textbf{Data Ingestion \& NLP}: Automated processing of Google Places metadata and text reviews, DistilBERT sentiment extraction, and keyword-masked ABSA feature calculation.
    \item \textbf{Core Modeling}: Machine learning models including SVD matrix factorization, content similarity matching, RRF hybrid fusion, and analytical aspect alignment scoring.
    \item \textbf{Conversational RAG}: Vector similarity search over 7,910 review chunks, prompt orchestration, local Ollama LLM inference (\texttt{mistral}/\texttt{llama3}), inline citation injection, and post-generation grounding validation.
    \item \textbf{Data Infrastructure}: Enterprise PostgreSQL 17.6 database server with \texttt{pgvector}, clean repository abstraction layers (`\texttt{PostgresHotelRepository}`, `\texttt{PgVectorEmbeddingRepository}`), SHA-256 field-level diff engine, FastAPI REST backend (`/api/v1/`), and a React + TypeScript frontend.
    \item \textbf{Orchestration \& Tooling}: Master CLI orchestrator, interactive ASCII progress monitoring dashboard (`\texttt{ProgressTracker}`), and structured run logging (`\texttt{pipeline.log}`).
\end{itemize}

Multi-city scaling outside Delhi NCR, multi-lingual NLP review parsing, and distributed database cluster deployments are outside the present implementation scope.

\section{Contributions}
This Independent Project delivers eight concrete system and engineering contributions:
\begin{enumerate}
    \item \textbf{Preference-Aware Hybrid Recommendation Engine}: Demonstrated that Reciprocal Rank Fusion (RRF) resolves score scale mismatches between collaborative and content-based models, elevating recommendation quality ($\text{NDCG}@10 > 0.12$).
    \item \textbf{Real-Time Analytical Explainability Engine}: Designed a deterministic feature-matching explainer that calculates aspect alignment scores without incurring the high latency of SHAP approximations.
    \item \textbf{Review-Grounded RAG Pipeline}: Combined hybrid vector retrieval with token-budget context compression and inline citation injection, backing LLM travel answers with explicit review text chunks.
    \item \textbf{Automated Grounding Validator}: Implemented a post-generation validation layer that cross-references LLM claims against retrieved evidence, stripping unverified amenity assertions ($96.7\%$ grounded response rate, $1.3\%$ hallucination rate).
    \item \textbf{Unified PostgreSQL + pgvector Migration}: Unified relational metadata and vector storage under PostgreSQL 17.6, verifying 1.0000 average embedding cosine similarity parity and zero orphan records across an 18-test provenance suite.
    \item \textbf{Repeatable Ingestion and Field-Level Diffing}: Implemented a 9-stage ingestion lifecycle utilizing SHA-256 canonical content hashing and selective vector synchronization to prevent redundant embedding recalculations.
    \item \textbf{One-Command CLI Orchestrator}: Developed a single-command CLI runner (`\texttt{python -m scripts.orchestrator full}`) unifying all 6 upstream processing stages.
    \item \textbf{Live Operational Visibility and Signal Safety}: Created an interactive ASCII terminal progress dashboard and a \texttt{SIGINT} (Ctrl+C) signal handler that terminates sub-processes cleanly, guaranteeing zero database corruption.
\end{enumerate}

\section{Organization of the Report}
The remainder of this report is organized into four subsequent chapters:
\begin{itemize}
    \item \textbf{Chapter 2: Literature Review} — Reviews foundational concepts and prior literature in hotel recommendation systems, collaborative filtering, aspect-based sentiment analysis, Retrieval-Augmented Generation, and vector database architectures.
    \item \textbf{Chapter 3: Methodology and System Architecture} — Provides a comprehensive technical breakdown of the system's 29-stage evolution, covering the data pipeline, NLP sentiment modeling, hybrid RRF recommendation logic, analytical explainability, grounded RAG architecture, PostgreSQL/pgvector database schema, and the repeatable CLI orchestration engine.
    \item \textbf{Chapter 4: Experimental Evaluation and Results} — Presents empirical evaluation results across recommender diagnostic benchmarks, retrieval ablation studies, grounding validator interception rates, pgvector backfill parity checks, and master backend provenance verification.
    \item \textbf{Chapter 5: Conclusion and Future Work} — Summarizes the primary engineering achievements of TrustLayer-AI and outlines future directions for multi-region scaling and automated vector synchronization.
\end{itemize}